\chapter{实验设计与结果分析}
\label{chap:results}
\section{实验设置}
本节应说明数据来源、样本划分、实验设备、软件环境与参数设置。对于涉及随机性的实验，应交代重复次数、随机种子和汇总方法，使读者能够理解结果产生的条件。模板不提供真实数据，也不预设任何方法优于其他方法。

对照方法的选择应与研究问题一致。使用相同的数据划分、计算预算和评价指标，必要时解释无法完全统一的条件。对于来自其他文献的设置，应给出对应来源，例如文献\cite{example-conference}。
\section{评价指标与比较结果}
先解释指标的含义与取值方向，再展示结果。分析时区分数值差异与实际意义，不应仅凭单次实验的最优值下结论。若使用\tabrefcn{tab:demo}的格式，应将所有虚构数据替换为可核查的真实结果。
\subsection{消融实验}
针对方法中的关键组成部分，设计能够单独检验其作用的实验。说明每个变体改变了什么、保持了什么，以及观察结果能支持哪些结论。
\subsection{局限与适用条件}
描述方法在哪些条件下有效、在哪些情况下性能下降，并结合证据讨论原因。对尚未验证的场景应如实说明，不要将推测写成已经证实的结论。
\section{本章小结}
按研究问题归纳实验结论，并明确证据能够支持的范围。避免重复逐项罗列图表中的所有数值。
